% ==============================================================================
% APÊNDICE B — Dataset do Experimento Controlado
% ==============================================================================

\chapter{Dataset do Experimento Controlado}
\label{ap:dataset}

Este apêndice apresenta o dataset completo dos 14 casos de teste utilizados no experimento controlado (OE6), incluindo expressões originais, gabaritos e resultados obtidos.

\section{Estrutura do Dataset}

O dataset foi armazenado em formato CSV com as seguintes colunas:

\begin{itemize}
    \item \textbf{id}: Identificador do caso de teste (ex: C1-01, C2-03)
    \item \textbf{cenario}: Tipo de mudança (Renomeio de Coluna, Tabela ou Medida)
    \item \textbf{descricao}: Descrição textual do caso de teste
    \item \textbf{change\_type}: Tipo de mudança (rename\_column, rename\_table, rename\_measure)
    \item \textbf{objeto}: Nome da medida ou coluna calculada impactada
    \item \textbf{old\_name}: Nome antigo do objeto renomeado
    \item \textbf{new\_name}: Nome novo do objeto renomeado
    \item \textbf{provedor}: Provedor LLM utilizado (groq, openai, anthropic)
    \item \textbf{modelo}: Modelo LLM específico (ex: llama-3.3-70b-versatile)
    \item \textbf{expressao\_original}: Expressão DAX original (entrada)
    \item \textbf{gabarito}: Expressão DAX esperada (gabarito)
    \item \textbf{expressao\_gerada}: Expressão DAX gerada pelo LLM (saída)
    \item \textbf{acuracia\_exata}: Booleano indicando se expressão gerada = gabarito
    \item \textbf{acuracia\_parcial}: Score de similaridade (0.0 a 1.0)
    \item \textbf{confianca\_llm}: Confiança estimada do LLM (0.0 a 1.0)
    \item \textbf{tempo\_s}: Tempo de resposta do LLM em segundos
    \item \textbf{erro}: Booleano indicando se houve erro de execução
    \item \textbf{erro\_msg}: Mensagem de erro (se aplicável)
\end{itemize}

\section{Casos de Teste Completos}

\subsection{Cenário 1: Renomeio de Coluna (7 casos)}

\subsubsection{Caso C1-01: Medida Simples com SUM}

\begin{itemize}
    \item \textbf{Descrição}: Medida simples com SUM referenciando a coluna renomeada
    \item \textbf{Mudança}: Renomear \texttt{Sales[Amount]} $\rightarrow$ \texttt{Sales[SalesAmount]}
    \item \textbf{Expressão Original}:
\end{itemize}

\begin{verbatim}
SUM(Sales[Amount])
\end{verbatim}

\begin{itemize}
    \item \textbf{Gabarito}:
\end{itemize}

\begin{verbatim}
SUM(Sales[SalesAmount])
\end{verbatim}

\begin{itemize}
    \item \textbf{Expressão Gerada}: \texttt{SUM(Sales[SalesAmount])}
    \item \textbf{Resultado}: ✅ Correto (acurácia = 100\%)
    \item \textbf{Tempo}: 1.04s
\end{itemize}

\subsubsection{Caso C1-02: CALCULATE com Filtro}

\begin{itemize}
    \item \textbf{Mudança}: Renomear \texttt{Sales[Amount]} $\rightarrow$ \texttt{Sales[SalesAmount]}
    \item \textbf{Expressão Original}:
\end{itemize}

\begin{verbatim}
CALCULATE(SUM(Sales[Amount]), Sales[Channel] = "Online")
\end{verbatim}

\begin{itemize}
    \item \textbf{Resultado}: ✅ Correto (acurácia = 100\%)
    \item \textbf{Tempo}: 0.81s
\end{itemize}

\subsubsection{Caso C1-03: SUMX Iterador}

\begin{itemize}
    \item \textbf{Mudança}: Renomear \texttt{Sales[Amount]} $\rightarrow$ \texttt{Sales[SalesAmount]}
    \item \textbf{Expressão Original}:
\end{itemize}

\begin{verbatim}
SUMX(Sales, Sales[Amount] * Sales[Quantity])
\end{verbatim}

\begin{itemize}
    \item \textbf{Resultado}: ✅ Correto (acurácia = 100\%)
    \item \textbf{Tempo}: 7.30s
\end{itemize}

\subsubsection{Caso C1-04: Coluna Calculada}

\begin{itemize}
    \item \textbf{Mudança}: Renomear \texttt{Sales[Amount]} $\rightarrow$ \texttt{Sales[SalesAmount]}
    \item \textbf{Expressão Original}:
\end{itemize}

\begin{verbatim}
Sales[Amount] * 1.1
\end{verbatim}

\begin{itemize}
    \item \textbf{Gabarito}: \texttt{Sales[SalesAmount] * 1.1}
    \item \textbf{Expressão Gerada}: \texttt{'Sales'[SalesAmount] * 1.1}
    \item \textbf{Resultado}: ❌ Incorreto (aspas desnecessárias)
    \item \textbf{Tempo}: 0.32s
\end{itemize}

\subsubsection{Caso C1-05: DIVIDE}

\begin{itemize}
    \item \textbf{Resultado}: ❌ Incorreto (aspas desnecessárias)
\end{itemize}

\subsubsection{Caso C1-06: Time Intelligence Complexo}

\begin{itemize}
    \item \textbf{Expressão Original}:
\end{itemize}

\begin{verbatim}
CALCULATE(
    SUM(Sales[Amount]),
    DATESYTD(Calendar[Date])
)
\end{verbatim}

\begin{itemize}
    \item \textbf{Resultado}: ✅ Correto (acurácia = 100\%)
    \item \textbf{Tempo}: 23.29s (expressão mais longa)
\end{itemize}

\subsubsection{Caso C1-07: Múltiplas Referências em VARs}

\begin{itemize}
    \item \textbf{Expressão Original}:
\end{itemize}

\begin{verbatim}
VAR CurrentSales = SUM(Sales[Amount])
VAR PrevSales = CALCULATE(SUM(Sales[Amount]), 
                          SAMEPERIODLASTYEAR(Calendar[Date]))
RETURN
    DIVIDE(CurrentSales - PrevSales, PrevSales)
\end{verbatim}

\begin{itemize}
    \item \textbf{Resultado}: ✅ Correto (acurácia = 100\%)
    \item \textbf{Tempo}: 0.36s
\end{itemize}

\subsection{Cenário 2: Renomeio de Tabela (4 casos)}

\subsubsection{Caso C2-01: SUM Simples}

\begin{itemize}
    \item \textbf{Mudança}: Renomear \texttt{Sales} $\rightarrow$ \texttt{SalesData}
    \item \textbf{Expressão Original}: \texttt{SUM(Sales[Amount])}
    \item \textbf{Gabarito}: \texttt{SUM(SalesData[Amount])}
    \item \textbf{Resultado}: ✅ Correto
    \item \textbf{Tempo}: 0.21s (caso mais rápido do experimento)
\end{itemize}

\subsubsection{Caso C2-02: CALCULATE + FILTER}

\begin{itemize}
    \item \textbf{Resultado}: ✅ Correto
\end{itemize}

\subsubsection{Caso C2-03: COUNTROWS}

\begin{itemize}
    \item \textbf{Expressão Original}: \texttt{COUNTROWS(Sales)}
    \item \textbf{Gabarito}: \texttt{COUNTROWS(SalesData)}
    \item \textbf{Expressão Gerada}: (resposta inválida)
    \item \textbf{Resultado}: ❌ Incorreto (LLM retornou resposta não-DAX)
    \item \textbf{Tempo}: 0.31s
    \item \textbf{Observação}: Expressão extremamente simples pode ter confundido o LLM
\end{itemize}

\subsubsection{Caso C2-04: SUMX + RELATED}

\begin{itemize}
    \item \textbf{Resultado}: ✅ Correto
    \item \textbf{Tempo}: 9.33s
\end{itemize}

\subsection{Cenário 3: Renomeio de Medida (3 casos)}

\subsubsection{Caso C3-01: Medida Simples}

\begin{itemize}
    \item \textbf{Mudança}: Renomear \texttt{[Total Sales]} $\rightarrow$ \texttt{[Total Revenue]}
    \item \textbf{Expressão Original}: \texttt{[Total Sales] - [Total Cost]}
    \item \textbf{Resultado}: ✅ Correto
    \item \textbf{Tempo}: 22.28s
\end{itemize}

\subsubsection{Caso C3-02: CALCULATE com Medida}

\begin{itemize}
    \item \textbf{Resultado}: ✅ Correto
\end{itemize}

\subsubsection{Caso C3-03: DIVIDE com Medidas}

\begin{itemize}
    \item \textbf{Resultado}: ✅ Correto
\end{itemize}

\section{Resumo Estatístico do Dataset}

\begin{table}[H]
\centering
\caption{Estatísticas do Dataset do Experimento}
\label{tab:dataset_stats}
\begin{tabular}{lc}
\toprule
\textbf{Métrica} & \textbf{Valor} \\
\midrule
Total de Casos & 14 \\
Casos Corretos & 11 (78.6\%) \\
Casos Incorretos & 3 (21.4\%) \\
Tempo Médio de Resposta & 4.96s \\
Tempo Mínimo & 0.21s (C2-01) \\
Tempo Máximo & 23.29s (C1-06) \\
Confiança Média do LLM & 92.86\% \\
Taxa de Erro do Sistema & 0.0\% \\
\bottomrule
\end{tabular}
\end{table}

\section{Disponibilidade do Dataset}

O dataset completo está disponível no repositório do projeto:

\begin{itemize}
    \item \textbf{CSV bruto}: \texttt{experiments/resultados/resultado\_20260505\_211431.csv}
    \item \textbf{Resumo JSON}: \texttt{experiments/resultados/resultado\_20260505\_211431\_resumo.json}
    \item \textbf{Análise}: \texttt{experiments/ANALISE\_RESULTADOS.md}
\end{itemize}

O repositório está publicado sob licença MIT em:

\url{https://github.com/isabellamoura/pbi-refactor-agent}

\section{Reprodutibilidade}

Para reproduzir o experimento:

\begin{verbatim}
# Clonar repositório
git clone https://github.com/isabellamoura/pbi-refactor-agent.git
cd pbi-refactor-agent

# Instalar dependências
pip install -e ".[dev]"

# Configurar API key (Groq é gratuito)
echo "GROQ_API_KEY=sua_chave_aqui" > .env

# Executar experimento
python experiments/experimento_controlado.py
\end{verbatim}

Os resultados serão salvos em \texttt{experiments/resultados/} com timestamp único para rastreabilidade.
